\section{Experimental design}
\label{sec:design}

\subsection{Setting and recruitment}
The study took place at a large public university in the Midwest. Employees can join the campus recreation center for 28 dollars a month, deducted from their pay. The center records each entry with an ID-card swipe. Members do not swipe out, so we count visits but cannot measure how long they last. \citet{acland2015}, by contrast, recorded both entry and exit swipes and could therefore measure how long each visit lasted. \citet{carrera2018incentives} introduced a swipe-out requirement partway through their study and credited a visit toward the incentive only after at least ten minutes in the gym.

In January 2025 we invited by email all 7,940 employees with an active membership to take part in a study of exercise habits. Of these, 2,537 gave informed consent and completed a baseline survey, for which they received a 10-dollar gift card. We excluded 125 employees who left the university or let their membership lapse before the start date, leaving 2,412 participants in 118 departments.

Enrollment and the baseline survey ran from 6 to 24 January 2025. We randomized departments on 27 January, and the intervention began on Monday 3 February 2025. The eight intervention weeks ended on 30 March, and the twelve follow-up weeks on 22 June. Weeks run from Monday to Sunday. A visit is a day with at least one entry swipe, so a second swipe on the same day does not count as a second visit, and the contract target of 12 visits refers to 12 different days.

\subsection{Randomization}
We randomized departments rather than individuals. Colleagues often exercise together and talk about workplace programs, and randomizing whole departments limits spillovers between arms. Randomization was stratified by department size (above or below 20 participants) and by the department's mean baseline attendance (above or below the median), which gives four strata. As Table~\ref{tab:balance} shows, 30, 29, 30 and 29 departments were assigned to the control, reminder, commitment and combined arms. All participants in a department received the same treatment.

\subsection{Treatments}
All arms received a one-page leaflet on the health benefits of regular exercise.

\emph{Reminders.} In the week before the intervention began, participants in the two reminder arms chose up to three weekdays for exercise on a short online form. During the eight intervention weeks they received a text message at 7 a.m. on each chosen day, for example: ``Today is one of your gym days. The recreation center opens at 6 a.m.''

\emph{Commitment contracts.} Participants in the two commitment arms received a contract offer by email in the first intervention week. The contract required at least 12 visits during the eight intervention weeks. Participants chose a stake between 20 and 100 dollars and paid it by card. The stake was refunded after week 8 if the participant met the target and was donated to a local food bank if not. Declining the offer had no consequences, and participants who signed could not change their stake later. The offer email explained the rules in four sentences and linked to a one-page form on which participants entered their stake. Participants who had not answered after four days received one follow-up email, and the offer closed at the end of the first intervention week. The staff of the recreation center did not know which participants held contracts.

\emph{Control.} The control group received the leaflet and nothing else.

No participant was paid for visiting the gym.

\subsection{Outcomes and registration}
We registered the design, the outcomes and the main specification before enrollment began. Like \citet{milkman2021}, who preregistered their megastudy as a single study before data collection, we registered a single primary outcome. The pre-registered primary outcome is the share of the eight contract weeks in which a participant visited the gym at least once; the number of visits per week is a secondary outcome. We measure persistence with the same two outcomes over weeks 9 to 20.

With 118 departments and about 600 participants per arm, a standard deviation of 0.9 visits per week and an intraclass correlation of 0.02 within departments, the design had 80 percent power at the 5 percent level for an MDE of 0.17 visits per week in any pairwise comparison of arms.

\subsection{Estimation}
We estimate intention-to-treat (ITT) effects with the regression
\begin{equation}
y_{it} = \alpha + \beta_1 C_i + \beta_2 R_i + \beta_3 C_i R_i + \gamma_t + \delta_{s(i)} + \varepsilon_{it},
\label{eq:main}
\end{equation}
where $y_{it}$ is an outcome of participant $i$ in week $t$, $C_i$ and $R_i$ indicate that the participant's department was assigned to the contract offer and to reminders, $\gamma_t$ are week fixed effects and $\delta_{s(i)}$ are fixed effects for the randomization stratum of the participant's department. The coefficient $\beta_1$ is the effect of the contract offer without reminders, $\beta_2$ the effect of reminders without the offer, and $\beta_3$ the change in the effect of the offer when reminders are added. Standard errors are clustered by participant because each participant contributes several weekly observations.
